% Options for packages loaded elsewhere
\PassOptionsToPackage{unicode}{hyperref}
\PassOptionsToPackage{hyphens}{url}
\documentclass[
  11pt,
]{article}
\usepackage{xcolor}
\usepackage[left=3cm,right=3cm,top=2.5cm,bottom=2.5cm]{geometry}
\usepackage{amsmath,amssymb}
\setcounter{secnumdepth}{5}
\usepackage{iftex}
\ifPDFTeX
  \usepackage[T1]{fontenc}
  \usepackage[utf8]{inputenc}
  \usepackage{textcomp} % provide euro and other symbols
\else % if luatex or xetex
  \usepackage{unicode-math} % this also loads fontspec
  \defaultfontfeatures{Scale=MatchLowercase}
  \defaultfontfeatures[\rmfamily]{Ligatures=TeX,Scale=1}
\fi
\usepackage{lmodern}
\ifPDFTeX\else
  % xetex/luatex font selection
\fi
% Use upquote if available, for straight quotes in verbatim environments
\IfFileExists{upquote.sty}{\usepackage{upquote}}{}
\IfFileExists{microtype.sty}{% use microtype if available
  \usepackage[]{microtype}
  \UseMicrotypeSet[protrusion]{basicmath} % disable protrusion for tt fonts
}{}
\usepackage{setspace}
\makeatletter
\@ifundefined{KOMAClassName}{% if non-KOMA class
  \IfFileExists{parskip.sty}{%
    \usepackage{parskip}
  }{% else
    \setlength{\parindent}{0pt}
    \setlength{\parskip}{6pt plus 2pt minus 1pt}}
}{% if KOMA class
  \KOMAoptions{parskip=half}}
\makeatother
\usepackage{longtable,booktabs,array}
\usepackage{caption}
\captionsetup[table]{skip=6pt}
\newcounter{none} % for unnumbered tables
\usepackage{calc} % for calculating minipage widths
% Correct order of tables after \paragraph or \subparagraph
\usepackage{etoolbox}
\makeatletter
\patchcmd\longtable{\par}{\if@noskipsec\mbox{}\fi\par}{}{}
\makeatother
% Allow footnotes in longtable head/foot
\IfFileExists{footnotehyper.sty}{\usepackage{footnotehyper}}{\usepackage{footnote}}
\makesavenoteenv{longtable}
\usepackage{graphicx}
\makeatletter
\newsavebox\pandoc@box
\newcommand*\pandocbounded[1]{% scales image to fit in text height/width
  \sbox\pandoc@box{#1}%
  \Gscale@div\@tempa{\textheight}{\dimexpr\ht\pandoc@box+\dp\pandoc@box\relax}%
  \Gscale@div\@tempb{\linewidth}{\wd\pandoc@box}%
  \ifdim\@tempb\p@<\@tempa\p@\let\@tempa\@tempb\fi% select the smaller of both
  \ifdim\@tempa\p@<\p@\scalebox{\@tempa}{\usebox\pandoc@box}%
  \else\usebox{\pandoc@box}%
  \fi%
}
% Set default figure placement to htbp
\def\fps@figure{htbp}
\makeatother
\setlength{\emergencystretch}{3em} % prevent overfull lines
\providecommand{\tightlist}{%
  \setlength{\itemsep}{0pt}\setlength{\parskip}{0pt}}
%% tools/stamp.tex
%% Provenance footer for PDFs rendered from this project.
%% Vendored by zzcollab; this file is not project-specific. The
%% footer prints the source document and the time of rendering.
%% The git version is defined here too, and so is available to a
%% project that wants it on the page, but it is left off the
%% default footer: it already names the staged copy in share/ and
%% appears in its manifest row. All values are supplied at render
%% time by tools/stamp-render.R, which writes a small generated
%% file that redefines the macros below. The \providecommand
%% defaults keep a bare LaTeX compile working when the document is
%% built without the wrapper.
\usepackage{fancyhdr}
\usepackage{xcolor}
\providecommand{\stampsource}{\detokenize{(source not recorded)}}
\providecommand{\stamptime}{(time not recorded)}
\providecommand{\stampversion}{\detokenize{(version not recorded)}}
\setlength{\footskip}{40pt}
\newcommand{\stampfootcontent}{%
  \footnotesize\color{gray}%
  \parbox[b]{\textwidth}{%
    \stamptime\hfill\thepage\\[2pt]%
    \ttfamily\raggedright\stampsource}}
\pagestyle{fancy}
\fancyhf{}
\fancyfoot[C]{\stampfootcontent}
\renewcommand{\headrulewidth}{0pt}
\renewcommand{\footrulewidth}{0.4pt}
%% The title page uses the `plain` page style; redefine it so the
%% stamp appears there too.
\fancypagestyle{plain}{%
  \fancyhf{}%
  \fancyfoot[C]{\stampfootcontent}%
  \renewcommand{\headrulewidth}{0pt}%
  \renewcommand{\footrulewidth}{0.4pt}}
\renewcommand{\stampsource}{\textasciitilde{}/\allowbreak{}prj/\allowbreak{}res/\allowbreak{}36-pmsimstats-ng/\allowbreak{}pmsimstats-ng/\allowbreak{}analysis/\allowbreak{}report/\allowbreak{}08-test-procedure-design-sensitivity/\allowbreak{}bullets.Rmd}
\renewcommand{\stamptime}{2026-08-15 16:51 PDT}
\renewcommand{\stampversion}{\detokenize{bc0b671-wip}}
\usepackage{amsmath}
\usepackage{amssymb}
\usepackage{booktabs}
\usepackage{longtable}
\usepackage{float}
\usepackage[mathlines]{lineno}
\linenumbers
\usepackage{bookmark}
\IfFileExists{xurl.sty}{\usepackage{xurl}}{} % add URL line breaks if available
\urlstyle{same}
% fallback for those not using the hyperref driver hyperxmp:
\makeatletter
\@ifundefined{xmpquote}{\newcommand{\xmpquote}[1]{#1}}{}
\makeatother
\hypersetup{
  pdftitle={Structured summary: Test-procedure choice for the biomarker-treatment interaction in aggregated N-of-1 trials: dichotomization{,} classical RM-ANOVA{,} mixed-effects{,} and generalized estimating equations},
  pdfauthor={pmsimstats team},
  hidelinks,
  pdfcreator={LaTeX via pandoc}}

\title{Structured summary: Test-procedure choice for the biomarker-treatment interaction in aggregated N-of-1 trials: dichotomization, classical RM-ANOVA, mixed-effects, and generalized estimating equations}
\author{pmsimstats team}
\date{August 15, 2026}

\begin{document}
\maketitle

{
\setcounter{tocdepth}{2}
\tableofcontents
}
\setstretch{1.1}
\emph{This document is a structured, section-by-section summary of the key points argued in the companion manuscript, ``Test-procedure choice for the biomarker-treatment interaction in aggregated N-of-1 trials: dichotomization, classical RM-ANOVA, mixed-effects, and generalized estimating equations.'' It is provided as a supplementary reading aid and does not stand on its own; all notation, derivations, simulation detail, and citations are in the main manuscript.}

\bigskip\noindent

\rule{\textwidth}{0.4pt}\bigskip

\subsection{A claim about how design and test choices interact}\label{a-claim-about-how-design-and-test-choices-interact}

\textbf{Joint optimization of test procedure and design parameters is the paper's central argument.}

\begin{itemize}\setlength{\itemsep}{2pt}\setlength{\parskip}{0pt}
\item The two principal choices -- test procedure and trial-design parameters -- are mutually dependent: validity and power of each test depend on the carryover and autocorrelation structure the design imposes.
\item Treating the choices as independent optimizations, as the published literature typically does, produces locally optimal but jointly suboptimal combinations.
\item The argument is presented as consequential at the sample sizes characteristic of precision-medicine pilot studies.
\end{itemize}

\textbf{The paper's two-part argument establishes a gap in the literature and then fills it at the test-procedure axis.}

\begin{itemize}\setlength{\itemsep}{2pt}\setlength{\parskip}{0pt}
\item No joint design-by-test sensitivity study exists for the aggregated N-of-1 design class.
\item The Monte Carlo experiment shows that linear-mixed with AR(1) dominates strict RM-ANOVA wherever residuals depart from compound symmetry.
\item Strict RM-ANOVA is retained as a conservative stress-test baseline rather than a primary recommendation.
\end{itemize}

\subsection{Why the question matters for chronic-condition pharmacology}\label{why-the-question-matters-for-chronic-condition-pharmacology}

\textbf{Physiological autocorrelation in chronic-condition trials violates the sphericity assumption that strict RM-ANOVA requires.}

\begin{itemize}\setlength{\itemsep}{2pt}\setlength{\parskip}{0pt}
\item Chronic-condition response trajectories exhibit serial correlation driven by drug pharmacokinetics, natural-history drift, and measurement reactivity.
\item Huynh-Feldt and Greenhouse-Geisser corrections recover nominal Type I only conditionally; the conventional threshold favors Greenhouse-Geisser when $\epsilon < 0.75$ (Blanca et al., 2023).
\item The mixed-effects model with \texttt{corCAR1} avoids these corrections by modeling AR(1) structure directly.
\end{itemize}

\textbf{Detecting the biomarker-treatment interaction is harder than detecting the average treatment effect (ATE) and is more sensitive to design choices.}

\begin{itemize}\setlength{\itemsep}{2pt}\setlength{\parskip}{0pt}
\item The PATH framework distinguishes effect-modeling from risk-modeling; the former requires more design-aware machinery.
\item Within-participant contrasts feeding the interaction estimator are sensitive to trial-design parameters in ways the ATE analysis is not.
\item Underpowered test procedures for the interaction will miss predictive-biomarker signal even at nominally adequate sample sizes.
\end{itemize}

\textbf{Cycle count and period length interact with carryover and test-procedure choice, so their joint optimum cannot be found by optimizing each separately.}

\begin{itemize}\setlength{\itemsep}{2pt}\setlength{\parskip}{0pt}
\item Additional cycles increase power for the biomarker-treatment interaction but with diminishing returns beyond the fourth or fifth cycle.
\item Shorter periods produce more cycles but worsen on-drug/off-drug separation when carryover is non-trivial.
\item The joint optimum on the (cycles, period length) plane depends on carryover half-life and test procedure together.
\end{itemize}

\subsection{What the published N-of-1 methodology literature does and does not do}\label{what-the-published-n-of-1-methodology-literature-does-and-does-not-do}

\textbf{The majority of published N-of-1 analyses are conducted under a misspecified residual covariance assumption.}

\begin{itemize}\setlength{\itemsep}{2pt}\setlength{\parskip}{0pt}
\item Systematic reviews document that 83.8\% of published N-of-1 studies ignore autocorrelation entirely.
\item CONSORT extension standards codify reporting quality but do not address the analytic-model decision.
\item The analytic-method gap has been identified as the highest-priority methodological problem in the field.
\end{itemize}

\textbf{Existing test-procedure comparisons address parallel-group or cluster designs but not the aggregated N-of-1 setting.}

\begin{itemize}\setlength{\itemsep}{2pt}\setlength{\parskip}{0pt}
\item The mixed model for repeated measures (MMRM) consensus applies to parallel-group longitudinal trials with monotone missingness, not N-of-1 designs.
\item Only 6.6\% of stepped-wedge cluster trials with fewer than six units use appropriate small-sample corrections.
\item No existing test-procedure comparison addresses the cross-over or N-of-1 setting directly.
\end{itemize}

\textbf{Design-parameter optimization has been studied in isolation from test-procedure choice, leaving the joint question unaddressed.}

\begin{itemize}\setlength{\itemsep}{2pt}\setlength{\parskip}{0pt}
\item Hendrickson et al. (2020) is the only paper comparing full design families for biomarker-interaction power, but holds period length fixed and uses a single test procedure.
\item Power calculations, autocorrelation effects, and cross-over design literature each address one axis without crossing them.
\item No published study crosses test procedure against the design-parameter grid under simulation for the aggregated N-of-1 class.
\end{itemize}

\textbf{No published study crosses test procedure against a design-parameter grid for the aggregated N-of-1 design class.}

\begin{itemize}\setlength{\itemsep}{2pt}\setlength{\parskip}{0pt}
\item The gap spans test-procedure comparison (strict RM-ANOVA, MMRM, GEE, mixed-effects with \texttt{corCAR1}) crossed against cycles, period lengths, asymmetry, run-in, and run-out.
\item The present paper addresses the gap on the test-procedure axis.
\item The design-parameter axis is deferred to a companion paper.
\end{itemize}

\subsection{Run-in design and the placebo-selection controversy}\label{run-in-design-and-the-placebo-selection-controversy}

\textbf{The Hendrickson hybrid design avoids placebo-run-in selection bias by using open-label stabilization, at some power cost.}

\begin{itemize}\setlength{\itemsep}{2pt}\setlength{\parskip}{0pt}
\item Response-based exclusion during a placebo run-in introduces selection bias that damages external and potentially internal validity (Berger et al., 2003).
\item The Hendrickson open-label stabilization phase avoids this bias but forgoes the power that aggressive response-based selection would supply.
\item The run-in choice is held fixed in the present paper; varying it is deferred to the companion design-sensitivity paper.
\end{itemize}

\subsection{What this paper contributes}\label{what-this-paper-contributes}

\textbf{The first contribution is a closed-form strict RM-ANOVA $F$-statistic for the biomarker-treatment interaction and its relation to the linear-mixed Wald test.}

\begin{itemize}\setlength{\itemsep}{2pt}\setlength{\parskip}{0pt}
\item The derivation specializes the balanced split-plot model to the biomarker-stratified case under the sphericity assumption.
\item The conditions under which the two tests agree and diverge are identified explicitly.
\item The derivation extends standard cross-over methodology to the biomarker-stratified case.
\end{itemize}

\textbf{The second contribution is specification of the cycle-by-period design grid calibrated to the prazosin/PTSD reference application.}

\begin{itemize}\setlength{\itemsep}{2pt}\setlength{\parskip}{0pt}
\item The grid varies cycle count, on-drug and off-drug period lengths, and on/off symmetry around the Hendrickson hybrid reference.
\item Calibration is to the prazosin/PTSD application that motivates the pmsimstats program.
\item The empirical sweep across this grid is deferred to a companion paper.
\end{itemize}

\textbf{The third contribution is a Monte Carlo simulation study crossing test procedures against the design grid under the ADEMP framework.}

\begin{itemize}\setlength{\itemsep}{2pt}\setlength{\parskip}{0pt}
\item The simulation runs on the pmsimstats package data-generating process under the mean-moderation architecture, exercised through \texttt{generateData()} and \texttt{lme\_analysis()}.
\item Five test procedures, including a continuous-biomarker contrast that isolates the dichotomization cost, are crossed against biomarker-moderation coefficient and sample size in a 40-cell factorial.
\item Results are reported following the Morris et al. (2019) ADEMP discipline including Monte Carlo standard errors (MCSE), bias, and interval coverage.
\end{itemize}

\textbf{The fourth contribution is a design-and-test recommendation table indexed by carryover half-life and target effect size.}

\begin{itemize}\setlength{\itemsep}{2pt}\setlength{\parskip}{0pt}
\item The recommendation is presented as preferred (cycle count, period length, test procedure) combinations rather than a single best design.
\item The recommendation is narrowed to the test-procedure axis at the reference design in the present paper.
\item The joint cycle-by-period optimization is reserved for the companion paper.
\end{itemize}

\textbf{The paper's organization maps each section to one of the four contributions.}

\begin{itemize}\setlength{\itemsep}{2pt}\setlength{\parskip}{0pt}
\item Sections 2 and 3 cover the closed-form derivation and the design-grid specification respectively.
\item Sections 4 and 5 present the simulation design and results.
\item Sections 6 and 7 synthesize the recommendation and discuss implications.
\end{itemize}

\section{Closed-form strict RM-ANOVA test of the biomarker-treatment interaction}\label{closed-form-strict-rm-anova-test-of-the-biomarker-treatment-interaction}

\textbf{The section derives the strict RM-ANOVA closed form and identifies three findings that anchor the rest of the paper.}

\begin{itemize}\setlength{\itemsep}{2pt}\setlength{\parskip}{0pt}
\item The derivation is confined to the strict classical compound-symmetric form; modern extensions appear only for context.
\item The simplest design supporting the test has two biomarker strata, two treatment phases, and $T$ within-phase timepoints.
\item Three findings from the derivation -- equivalence to a two-sample $t$-test, a closed-form power formula, and five data-generating process (DGP) violations -- anchor the simulation study.
\end{itemize}

\subsection{Notation and design}\label{notation-and-design}

\textbf{The notation section establishes a restricted design that is deliberately simpler than the pmsimstats DGP.}

\begin{itemize}\setlength{\itemsep}{2pt}\setlength{\parskip}{0pt}
\item The design is the simplest that supports a biomarker-treatment interaction test in a within-subject setting.
\item The gap between this restricted design and the pmsimstats DGP is addressed explicitly in Section 2.4.
\item All departures from the DGP are flagged as sources of conservative bias in the strict RM-ANOVA operating characteristics.
\end{itemize}

\textbf{The split-plot structure arises because the biomarker stratum is between-subjects while treatment phase and timepoints are within-subjects.}

\begin{itemize}\setlength{\itemsep}{2pt}\setlength{\parskip}{0pt}
\item Strict RM-ANOVA requires categorical predictors; dichotomizing the continuous biomarker at a cut-point is the price of admission.
\item Each subject is nested in exactly one biomarker stratum but passes through both treatment phases.
\item The design is therefore a canonical split-plot repeated-measures design.
\end{itemize}

\textbf{The classical split-plot model encodes the biomarker-treatment interaction as the $(\alpha\beta)_{ij}$ fixed effect under compound symmetry.}

\begin{itemize}\setlength{\itemsep}{2pt}\setlength{\parskip}{0pt}
\item Sum-to-zero constraints are imposed on all fixed effects; the random subject effect is nested in biomarker stratum.
\item The sphericity (compound symmetry) condition requires equal within-subject covariance across all observation pairs.
\item The biomarker-treatment interaction $(\alpha\beta)_{ij}$ and its time-resolved form $(\alpha\beta\tau)_{ijk}$ are the primary targets.
\end{itemize}

\subsection{\texorpdfstring{Sums of squares and the interaction \(F\)-statistic}{Sums of squares and the interaction F-statistic}}\label{sums-of-squares-and-the-interaction-f-statistic}

\textbf{The total sum of squares partitions into between-subjects and within-subjects components that are orthogonal under balance.}

\begin{itemize}\setlength{\itemsep}{2pt}\setlength{\parskip}{0pt}
\item The between-subjects component captures biomarker-stratum effects and subject-within-stratum variation.
\item The within-subjects component captures treatment-phase effects, the biomarker-treatment interaction, and within-subjects error.
\item Orthogonality under balance ensures the interaction $F$-ratio uses the correct within-subjects error term.
\end{itemize}

\textbf{The interaction $F$-ratio uses the within-subjects error $SS_{B \times S(A)}$ as its denominator, with degrees of freedom $a(n-1)(b-1)$.}

\begin{itemize}\setlength{\itemsep}{2pt}\setlength{\parskip}{0pt}
\item $SS_{AB}$ has $df_{AB} = (a-1)(b-1)$ degrees of freedom; the within-subjects error has $df_{B \times S(A)} = a(n-1)(b-1)$.
\item The within-subjects error term is the appropriate denominator for the treatment effect and its interaction with the between-subjects factor.
\item The $F$-ratio $F_{AB} = MS_{AB}/MS_{B \times S(A)}$ follows a central $F$ under $H_0$ and sphericity.
\end{itemize}

\subsubsection{\texorpdfstring{The \(2 \times 2\) split-plot case}{The 2 \textbackslash times 2 split-plot case}}\label{the-2-times-2-split-plot-case}

\textbf{Specializing to $a = b = 2$ with $T = 1$ reduces the interaction to the difference of within-subject treatment differences across biomarker strata.}

\begin{itemize}\setlength{\itemsep}{2pt}\setlength{\parskip}{0pt}
\item Setting $T = 1$ corresponds to pre-averaging within-phase observations before computing the interaction.
\item The sample interaction contrast $\hat{L}$ is the difference of within-subject treatment differences across strata.
\item This specialization yields the cleanest closed-form expression for the interaction $F$-statistic.
\end{itemize}

\textbf{Under sphericity and balance in the $2 \times 2$ case, the interaction $F$-statistic reduces to $nT\hat{L}^2 / (4\hat{\sigma}_e^2)$ with degrees of freedom $(1, 2(n-1))$.}

\begin{itemize}\setlength{\itemsep}{2pt}\setlength{\parskip}{0pt}
\item The interaction sum of squares is $SS_{AB} = nT\hat{L}^2/4$; with $df_{AB} = 1$ this equals the interaction mean square.
\item Under sphericity, $MS_{B \times S(A)} = \hat{\sigma}_e^2$, yielding the closed form directly.
\item The error degrees of freedom are $2(n-1)$, reflecting one degree of freedom per subject in each of the two strata.
\end{itemize}

\textbf{Equation \ref{eq:rmanova-f-22} is the paper's strict RM-ANOVA $F$ reference expression and is named to distinguish it from the mixed-effects and GEE Wald statistics.}

\begin{itemize}\setlength{\itemsep}{2pt}\setlength{\parskip}{0pt}
\item The expression applies to the $2 \times 2 \times T$ split-plot design and is the textbook compound-symmetric form.
\item It serves as the analytic anchor for the power formula in Section 2.3.1 and for comparison with simulation results in Section 5.
\item The name 'strict RM-ANOVA $F$' is used throughout to mark the distinction from Wald statistics that do not require sphericity.
\end{itemize}

\subsubsection{\texorpdfstring{Equivalence to a within-subject \(t\)-test}{Equivalence to a within-subject t-test}}\label{equivalence-to-a-within-subject-t-test}

\textbf{The strict RM-ANOVA $F$-test for the biomarker-treatment interaction is equivalent to a two-sample $t$-test on within-subject treatment differences.}

\begin{itemize}\setlength{\itemsep}{2pt}\setlength{\parskip}{0pt}
\item The within-subject treatment difference $\Delta_{il}$ cancels the random subject effect, leaving a pure within-subject contrast.
\item $\hat{L} = \bar{\Delta}_1 - \bar{\Delta}_2$ is the difference of stratum-mean treatment differences.
\item The equivalence $t_{AB}^2 = F_{AB}$ is well-known in the split-plot literature and makes the test's assumptions unusually transparent.
\end{itemize}

\textbf{$t_{AB}^2 = F_{AB}$ recovers the interaction $F$-statistic, confirming that the strict RM-ANOVA test is a two-sample comparison of within-subject treatment effects.}

\begin{itemize}\setlength{\itemsep}{2pt}\setlength{\parskip}{0pt}
\item The identity holds under sphericity because the random subject effect cancels in within-subject differences.
\item The test reduces to asking whether the within-subject treatment effect differs between biomarker strata.
\item Transparency about this reduction is analytically useful: anything that biases within-subject differences distorts the test in a predictable direction.
\end{itemize}

\textbf{The transparent reduction to a two-sample $t$-test will be used in Section 5 to interpret the simulation results.}

\begin{itemize}\setlength{\itemsep}{2pt}\setlength{\parskip}{0pt}
\item The reduction is exploited in Section 5 to connect strict RM-ANOVA results to the mixed-effects findings.
\item Biases and inflations in within-subject differences produce predictable distortions in the test's operating characteristics.
\item The intuitive appeal of the reduction makes the test's limitations equally transparent.
\end{itemize}

\textbf{The general $a \times b$ case and the time-resolved three-way interaction extend equation \ref{eq:rmanova-f} with degrees of freedom scaled by $(a-1)(b-1)$.}

\begin{itemize}\setlength{\itemsep}{2pt}\setlength{\parskip}{0pt}
\item For arbitrary $a$ and $b$, the interaction $F$-statistic uses $df_1 = (a-1)(b-1)$ and $df_2 = a(n-1)(b-1)$.
\item The three-way interaction $(\alpha\beta\tau)_{ijk}$ tests whether the biomarker modifies the treatment effect across time.
\item The corresponding degrees of freedom are $(a-1)(b-1)(T-1)$ and $a(n-1)(b-1)(T-1)$.
\end{itemize}

\subsection{Relation to the linear-mixed Wald test}\label{relation-to-the-linear-mixed-wald-test}

\textbf{The strict RM-ANOVA $F$ and the linear-mixed Wald $t$ are asymptotically equivalent only under three conditions that the pmsimstats DGP does not satisfy simultaneously.}

\begin{itemize}\setlength{\itemsep}{2pt}\setlength{\parskip}{0pt}
\item Equivalence requires dichotomized biomarker, compound-symmetric within-subject covariance, and binary \texttt{Dbc}.
\item When any condition fails, the mixed-model test extracts more information via continuous biomarker, AR(1) modeling, or continuous exposure-decay.
\item In the pmsimstats setting, all three conditions fail, so the mixed-model test dominates strictly.
\end{itemize}

\textbf{Sphericity violations inflate Type I error; Greenhouse-Geisser and Huynh-Feldt corrections recover nominal control by penalising degrees of freedom.}

\begin{itemize}\setlength{\itemsep}{2pt}\setlength{\parskip}{0pt}
\item Mauchly's test detects compound-symmetry failure but is not informative about direction; Box's $\epsilon$ is a scalar summary with $\epsilon = 1$ at perfect sphericity.
\item The Greenhouse-Geisser correction adjusts both numerator and denominator degrees of freedom by $\hat{\epsilon}$; Huynh-Feldt is less conservative and is recommended when $\hat{\epsilon} > 0.75$.
\item The conventional rule recommends Greenhouse-Geisser when $\hat{\epsilon} < 0.75$ to avoid substantial Type I error inflation.
\end{itemize}

\textbf{Under the pmsimstats AR(1) DGP, a mixed model with no residual-correlation structure inflates Type I error to 13--17\%; \texttt{corCAR1} restores nominal control, and the pre-averaged strict RM-ANOVA arm is sphericity-robust.}

\begin{itemize}\setlength{\itemsep}{2pt}\setlength{\parskip}{0pt}
\item AR(1) within-factor correlation is not compound-symmetric; Box's $\epsilon$ is strictly below 1.
\item Greenhouse-Geisser correction recovers nominal Type I error but substantially attenuates power by penalising degrees of freedom.
\item The \texttt{corCAR1} mixed-model analysis models AR(1) structure directly and does not pay the degrees-of-freedom penalty.
\end{itemize}

\textbf{The strict RM-ANOVA test is a conservative reduction rather than a competing alternative; option (c) -- reporting both with explicit power-loss tabulation -- is operationalized in the simulation study.}

\begin{itemize}\setlength{\itemsep}{2pt}\setlength{\parskip}{0pt}
\item Three defensible responses to a reviewer requesting RM-ANOVA results are identified, ranging from reframing the LME analysis to reporting both with power-loss tables.
\item Option (c), reporting both procedures with explicit power-loss tabulation, is the strategy adopted in Sections 4--5.
\item This strategy is the most transparent and the most informative for the reviewer.
\end{itemize}

\subsubsection{Power calculation under the closed form}\label{power-calculation-under-the-closed-form}

\textbf{The non-centrality parameter $\lambda = nT\delta^2 / (4\sigma_e^2)$ provides the analytic power anchor for the simulation study.}

\begin{itemize}\setlength{\itemsep}{2pt}\setlength{\parskip}{0pt}
\item Under the alternative $H_1: \hat{L} = \delta$, the $F$-statistic follows a non-central $F$ distribution.
\item The non-centrality parameter scales linearly with $n$ and $T$ and quadratically with $\delta$.
\item The formula inverts to a sample-size formula $n = 4\sigma_e^2\lambda / (T\delta^2)$ for power planning.
\end{itemize}

\textbf{All five assumptions of the closed-form power formula are violated to varying degrees in the pmsimstats DGP, so the formula provides an upper bound rather than a calibrated prediction.}

\begin{itemize}\setlength{\itemsep}{2pt}\setlength{\parskip}{0pt}
\item The five assumptions are: sphericity, categorical biomarker with equal stratum sizes, no carryover, no within-treatment relapse, and i.i.d. within-phase observations.
\item Violations of sphericity and i.i.d. within-phase observations are two faces of the same AR(1) departure from compound symmetry.
\item The analytic power formula is best understood as an upper bound on power available from strict RM-ANOVA under the pmsimstats DGP.
\end{itemize}

\textbf{The five assumptions are not independent: sphericity and i.i.d. within-phase observations are two faces of the same AR(1) departure from compound symmetry.}

\begin{itemize}\setlength{\itemsep}{2pt}\setlength{\parskip}{0pt}
\item Under AR(1), compound-symmetry failure drives both Type I error inflation and within-phase autocorrelation simultaneously.
\item The analytic power formula is therefore a joint upper bound, not a calibrated predictor for either failure mode in isolation.
\item This dependence motivates the Monte Carlo validation in Sections 4--5 rather than reliance on the closed form alone.
\end{itemize}

\subsubsection{Limitations of the strict classical formulation}\label{limitations-of-the-strict-classical-formulation}

\textbf{The five assumptions of the strict classical formulation are empirically problematic in the pmsimstats setting; their severity is presented in order.}

\begin{itemize}\setlength{\itemsep}{2pt}\setlength{\parskip}{0pt}
\item Closed-form simplicity is bought at the price of categorical biomarker, sphericity, no carryover model, balanced design, and equal time spacing.
\item Each assumption is violated by the pmsimstats DGP to a degree that favors the linear-mixed analysis.
\item The most severe violations -- categorical biomarker and no carryover -- together account for a 30--50\% attenuation of the interaction estimator.
\end{itemize}

\textbf{The five limitations motivate the primary-and-sensitivity reporting strategy; the categorical-biomarker and no-carryover violations are the most consequential.}

\begin{itemize}\setlength{\itemsep}{2pt}\setlength{\parskip}{0pt}
\item The categorical-biomarker and no-carryover violations together account for a 30--50\% attenuation of the interaction estimator relative to the linear-mixed analysis.
\item The linear-mixed analysis with \texttt{corCAR1} is retained as primary; the strict RM-ANOVA result is reported as a sensitivity comparator.
\item The Section 6 recommendation operationalizes this strategy for typical chronic-condition N-of-1 scenarios.
\end{itemize}

\section{Cycle-by-period design grid}\label{cycle-by-period-design-grid}

\textbf{The cycle-by-period grid is specified here but its empirical sweep is deferred; the Section 6 recommendation is correspondingly narrowed to the test-procedure axis.}

\begin{itemize}\setlength{\itemsep}{2pt}\setlength{\parskip}{0pt}
\item The design grid varies cycle count, on-drug and off-drug period lengths, and on/off symmetry around the Hendrickson hybrid reference.
\item The present paper's scope is narrowed to the test-procedure axis at the prazosin/PTSD reference design.
\item The joint (cycle count, period length, test procedure) optimum is deferred to the follow-up paper and recorded as a portfolio item.
\end{itemize}

\section{Simulation design}\label{simulation-design}

\textbf{The ADEMP framework structures the simulation program around the biomarker-treatment interaction estimand under five test procedures.}

\begin{itemize}\setlength{\itemsep}{2pt}\setlength{\parskip}{0pt}
\item The five procedures are strict RM-ANOVA (dichotomized biomarker), a continuous-biomarker within-subject contrast, linear-mixed with \texttt{corCAR1}, GEE with naive sandwich, and GEE with Mancl-DeRouen small-sample correction.
\item The continuous-biomarker contrast separates the cost of dichotomization from the cost of the test procedure proper; the four continuous-biomarker arms share a common interaction estimand.
\item Primary outputs are power and Type I error at $\alpha = 0.05$; secondary estimands (bias, empirical and model SE, 95\% CI coverage) are reported for the four coefficient-bearing arms.
\item All estimates are reported with MCSE following the Morris et al. (2019) discipline.
\end{itemize}

\textbf{MCSE for power and Type I error follow the binomial-proportion form and are reported alongside every estimate.}

\begin{itemize}\setlength{\itemsep}{2pt}\setlength{\parskip}{0pt}
\item Power and Type I error MCSE use $\sqrt{\hat{\pi}(1-\hat{\pi})/n_{\text{reps}}}$.
\item Bias and empirical SE of $\hat{\beta}_{bm:D}$ use the standard within-replicate mean estimator.
\item Reporting MCSE alongside every estimate is the Morris et al. discipline adopted project-wide.
\end{itemize}

\textbf{The full ADEMP pre-registration fixes three studies before production runs commit, providing a pre-specified analytic plan.}

\begin{itemize}\setlength{\itemsep}{2pt}\setlength{\parskip}{0pt}
\item Study 1 is the test-procedure comparison at fixed design (27 cells); Studies 2 and 3 cover the cycle-by-period grid and joint heatmap (240 cells each).
\item Pre-registration at \texttt{00-ademp-pre-registration.md} locks the factorial before production runs commit.
\item The present paper reports only Study 1; Studies 2 and 3 are deferred to the companion paper.
\end{itemize}

\textbf{The factorial crosses five procedures by two $\beta_{bm}$ values by four $N$ values; the inline Mancl-DeRouen SE is benchmarked against \texttt{geesmv::GEE.var.md}.}

\begin{itemize}\setlength{\itemsep}{2pt}\setlength{\parskip}{0pt}
\item The simulation crosses five procedures by $\beta_{bm} \in \{0, 0.45\}$ by $N \in \{25, 35, 70, 100\}$, with 1000 replicates per power cell and 2000 per Type I cell.
\item The Mancl-DeRouen correction inflates each cluster's residual via $(I - H_{ii})^{-1}r_i$ in the sandwich meat.
\item The inline SE agrees with \texttt{geesmv::GEE.var.md} to within roughly 2.5\% (mean ratio 1.02 over 25 probe replicates), so the correction magnitude is externally validated.
\end{itemize}

\section{Results}\label{results}

\textbf{The 40-cell factorial completed in 587 s with 100\% convergence across 60{,}000 fits; MCSE is approximately 0.016 for power and 0.005 for Type I error.}

\begin{itemize}\setlength{\itemsep}{2pt}\setlength{\parskip}{0pt}
\item Execution used 7-core \texttt{parallel::mclapply} with L'Ecuyer-CMRG streams; convergence was 100\% in every cell.
\item Per-replicate output and the ADEMP summary are archived at \texttt{analysis/data/quick-sim/}.
\item MCSE on power near 0.5 is 0.016 (1000 reps); MCSE on Type I error near 0.05 is 0.005 (2000 reps), providing adequate resolution to distinguish the procedures.
\end{itemize}

\subsection{Test-procedure comparison at fixed design}\label{test-procedure-comparison-at-fixed-design}

\textbf{Strict RM-ANOVA targets a dichotomized estimand; the continuous-biomarker contrast isolates the cost of dichotomization from the cost of the test procedure proper.}

\begin{itemize}\setlength{\itemsep}{2pt}\setlength{\parskip}{0pt}
\item Strict RM-ANOVA operates on a median-dichotomized biomarker, pre-averaged to one value per phase; its estimand is on the dichotomized scale.
\item The continuous-biomarker contrast uses the same pre-averaging but the continuous biomarker, so it shares the estimand of the linear-mixed and GEE arms.
\item The three-rung ladder (RM-ANOVA, contrast, linear-mixed) decomposes the gap between classical RM-ANOVA and the mixed model into a dichotomization component and a test-procedure component.
\end{itemize}

\textbf{Most of the gap between strict RM-ANOVA and the mixed model is the cost of dichotomization, not of the test procedure: the continuous-biomarker contrast recovers four-fifths of it.}

\begin{itemize}\setlength{\itemsep}{2pt}\setlength{\parskip}{0pt}
\item At $N = 25$, replacing the dichotomized biomarker with the continuous one (RM-ANOVA $\to$ contrast) raises power from 0.358 to 0.538, a gain of 0.180.
\item Moving from the pre-averaged contrast to the full linear-mixed model adds only 0.027 (0.538 $\to$ 0.565); the test-procedure component is small.
\item The same pattern holds across $N$: dichotomization costs 0.16--0.21 of power, while the longitudinal test procedure adds at most 0.04.
\end{itemize}

\textbf{Naive GEE inflates Type I error across the whole sample-size range, not only at small $N$; the Mancl-DeRouen correction restores nominal control at a moderate power cost.}

\begin{itemize}\setlength{\itemsep}{2pt}\setlength{\parskip}{0pt}
\item Naive GEE rejects under the null at 0.097, 0.086, 0.079, and 0.074 for $N = 25, 35, 70, 100$ -- between $1.5\times$ and $1.9\times$ nominal throughout.
\item The Mancl-DeRouen correction brings the rate to 0.064, 0.053, 0.063, and 0.059, within or close to MCSE of the nominal 0.05 at every $N$.
\item The correction's power cost is moderate (about 0.10 at $N = 25$, 0.07 at $N = 35$) and the corrected procedure no longer over-rejects.
\end{itemize}

\textbf{Standard-error calibration separates the four coefficient-bearing arms: the contrast and Mancl-DeRouen GEE are well-calibrated, naive GEE under-covers, and the linear-mixed model over-covers.}

\begin{itemize}\setlength{\itemsep}{2pt}\setlength{\parskip}{0pt}
\item The continuous-biomarker contrast and Mancl-DeRouen GEE attain 93--95\% coverage with model/empirical SE ratios near 1.
\item Naive GEE under-covers at small $N$ (90.0\% at $N = 25$) because its model SE is 9--12\% too small.
\item The linear-mixed model over-covers (97.0--97.7\%) because its model SE is 11--16\% larger than the empirical SE, the calibration signature of its conservatism.
\end{itemize}

\textbf{All four coefficient-bearing arms are approximately unbiased; the linear-mixed conservatism is a variance-calibration effect, not a point-estimate effect, and a Kenward-Roger correction would address it.}

\begin{itemize}\setlength{\itemsep}{2pt}\setlength{\parskip}{0pt}
\item Across all cells, the relative bias of the interaction coefficient is below 4\% of the estimand for every coefficient-bearing arm.
\item The linear-mixed Type I error is 0.021 to 0.034, below nominal, consistent with \texttt{corCAR1} degrees-of-freedom accounting under $T = 8$ timepoints.
\item A Kenward-Roger or Satterthwaite small-sample correction would address the conservatism and is recorded as a follow-up item.
\end{itemize}

\section{Design-and-test recommendation}\label{design-and-test-recommendation}

\textbf{The recommendation is presented as a table of preferred (procedure, scenario) combinations at the reference design rather than as a single best procedure.}

\begin{itemize}\setlength{\itemsep}{2pt}\setlength{\parskip}{0pt}
\item Cycle count and period length are held at reference values; joint optimization over those axes is deferred to the follow-up paper.
\item The preferred procedure depends on features of the application -- sample size, tolerance for Type I inflation, and reviewer requirements -- that the analyst must determine.
\item The first-order recommendation is to keep the biomarker continuous; the choice among continuous-biomarker procedures is second-order.
\end{itemize}

\textbf{Keep the biomarker continuous; among continuous-biomarker procedures the linear-mixed model is the preferred default, and naive GEE should not be used without a small-sample correction at any $N$ examined.}

\begin{itemize}\setlength{\itemsep}{2pt}\setlength{\parskip}{0pt}
\item The first-order decision is to retain the continuous biomarker; dichotomization forfeits 0.08--0.21 of power across the grid.
\item Among continuous-biomarker procedures, linear-mixed with \texttt{corCAR1} has the highest power within nominal Type I; the contrast is a close, transparent alternative.
\item Naive GEE is anti-conservative at every $N$ examined and should be replaced by the Mancl-DeRouen correction where a sandwich estimator is required.
\end{itemize}

\subsection{What the Monte Carlo study settles}\label{what-the-monte-carlo-study-settles}

\textbf{The simulation settles four substantive questions; the inline Mancl-DeRouen SE is validated against \texttt{geesmv::GEE.var.md} to within 2.5\%.}

\begin{itemize}\setlength{\itemsep}{2pt}\setlength{\parskip}{0pt}
\item The 40-cell factorial completed in 587 s with 100\% convergence across 60{,}000 fits; MCSE is 0.016 for power near 0.5 and 0.005 for Type I error near 0.05.
\item The inline Mancl-DeRouen implementation produces identical point estimates and larger standard errors than the naive sandwich.
\item Its SE for the interaction coefficient matches \texttt{geesmv::GEE.var.md} to within roughly 2.5\% over 25 probe replicates, so the correction magnitude is externally validated.
\end{itemize}

\textbf{The first resolved question is that the classical-versus-mixed power gap is mostly the cost of dichotomization: the continuous-biomarker contrast recovers about four-fifths of it.}

\begin{itemize}\setlength{\itemsep}{2pt}\setlength{\parskip}{0pt}
\item At $N = 25$, dichotomization (RM-ANOVA $\to$ contrast) costs 0.180 of power; the test procedure (contrast $\to$ linear-mixed) adds 0.027.
\item The decomposition is stable across $N$: dichotomization costs 0.08--0.21, the test procedure at most 0.04.
\item A classical, sphericity-free, AR(1)-free contrast on the continuous biomarker therefore captures most of the available power.
\end{itemize}

\textbf{The second resolved question is that naive GEE inflates Type I error across the entire sample-size range, and Mancl-DeRouen restores nominal control.}

\begin{itemize}\setlength{\itemsep}{2pt}\setlength{\parskip}{0pt}
\item Naive GEE rejects at 0.097, 0.086, 0.079, 0.074 for $N = 25, 35, 70, 100$ -- still about $1.5\times$ nominal at $N = 100$.
\item The Mancl-DeRouen correction brings the rate to 0.064, 0.053, 0.063, 0.059, within or near MCSE of nominal at every $N$.
\item The power cost is moderate (about 0.10 at $N = 25$, 0.07 at $N = 35$).
\end{itemize}

\textbf{The third resolved question is standard-error calibration: the contrast and Mancl-DeRouen GEE are well-calibrated, naive GEE under-covers, and the linear-mixed model over-covers.}

\begin{itemize}\setlength{\itemsep}{2pt}\setlength{\parskip}{0pt}
\item Contrast and Mancl-DeRouen GEE attain 93--95\% coverage with model/empirical SE ratios near 1.
\item Naive GEE under-covers (90.0\% at $N = 25$) because its model SE is 9--12\% too small.
\item The linear-mixed model over-covers (97.0--97.7\%) because its model SE exceeds the empirical SE by 11--16\%.
\end{itemize}

\textbf{The fourth resolved question is that the linear-mixed conservatism is a degrees-of-freedom artifact a Kenward-Roger correction would remove.}

\begin{itemize}\setlength{\itemsep}{2pt}\setlength{\parskip}{0pt}
\item LME Type I error is 0.021 to 0.034, conservative rather than nominal, consistent with \texttt{corCAR1} degrees-of-freedom accounting under $T = 8$ timepoints.
\item The conservatism shows up as over-coverage and a model/empirical SE ratio above one, not as point-estimate bias.
\item A Kenward-Roger or Satterthwaite correction would recover the small power increment the conservatism forgoes.
\end{itemize}

\section{Conclusions}\label{conclusions}

\textbf{The biomarker should be kept continuous: dichotomization, not the test procedure, drives most of the classical-versus-mixed power gap; naive GEE is disqualified without a small-sample correction.}

\begin{itemize}\setlength{\itemsep}{2pt}\setlength{\parskip}{0pt}
\item Dichotomizing the biomarker for strict RM-ANOVA forfeits 0.08--0.21 of power across $N$; the test procedure proper adds at most 0.04.
\item Among continuous-biomarker procedures the linear-mixed model has the highest power within nominal Type I; the contrast is a close, transparent alternative.
\item Naive GEE inflates Type I error to 1.5--1.9$\times$ nominal across all $N$; the Mancl-DeRouen correction restores nominal control at a moderate (about 10-percentage-point) power cost.
\end{itemize}

\textbf{The closed-form reduction to a two-sample $t$-test makes the strict RM-ANOVA transparent; in this study its limitation is dichotomization, with sphericity relevant only to the un-pre-averaged classical analysis.}

\begin{itemize}\setlength{\itemsep}{2pt}\setlength{\parskip}{0pt}
\item Under sphericity and balance, the biomarker-treatment interaction $F$-test reduces to a two-sample $t$-test on within-subject treatment differences.
\item The pre-averaged-to-one-value-per-phase form evaluated here is sphericity-robust by construction and holds nominal Type I error, so its power loss is attributable to dichotomization, not to sphericity.
\item Sphericity violation under AR(1) bites only when the full multi-timepoint series is analyzed by classical RM-ANOVA without a degrees-of-freedom correction; the linear-mixed model with \texttt{corCAR1} sidesteps both issues.
\end{itemize}

\textbf{The cycle-by-period design grid is deferred to a follow-up paper; infrastructure and design specification are in place and the scope limitation is explicitly acknowledged.}

\begin{itemize}\setlength{\itemsep}{2pt}\setlength{\parskip}{0pt}
\item The simulation infrastructure for the cycle-by-period sweep is committed at \texttt{analysis/scripts/test-procedure-design-sensitivity/}.
\item The design specification is at \texttt{docs/27-cycle-period-design-sweep-plan.md}.
\item The follow-up paper will deliver the joint recommendation table that the present paper's Section 6 narrows to the test-procedure-only axis.
\end{itemize}

\section{Data availability statement}\label{data-availability-statement}

\textbf{Data and code are openly available in the pmsimstats-ng repository at paths listed in the Reproducibility section.}

\begin{itemize}\setlength{\itemsep}{2pt}\setlength{\parskip}{0pt}
\item Per-replicate Monte Carlo outputs and ADEMP cell-level summaries are versioned under \texttt{analysis/data/}.
\item Simulation drivers are at \texttt{analysis/scripts/}.
\item Exact paths relevant to this manuscript are listed in the Reproducibility section.
\end{itemize}

\end{document}
